\section{Exercices d'application}

\rubrique{Positions relatives de droites et plans}

\exo{}\diff{1}
On considère un cube $ABCDEFGH$. 
\begin{enumerate}[label=\alph*)]
\item Donner la position relative de la droite $(AB)$ et du plan $(EFG)$.
\item Donner la position relative de la droite $(AC)$ et du plan $(EFH)$.
\item Donner la position relative des droites $(AB)$ et $(GH)$.
\item Donner la position relative des droites $(AC)$ et $(FH)$.
\end{enumerate}

\exo{}\diff{1}
Soit $ABCD$ un tétraèdre. $I$ est un point de l'arête $[AB]$, $J$ un point de $[AC]$ et $K$ un point de $[AD]$. 
\begin{enumerate}[label=\alph*)]
\item Que peut-on dire des droites $(IJ)$ et $(BC)$ ? Justifier.
\item Que peut-on dire des plans $(IJK)$ et $(BCD)$ ? Justifier.
\end{enumerate}

\exo{}\diff{2}
Dans un repère orthonormé de l'espace $(O;\vec{i},\vec{j},\vec{k})$, on considère les points $A(1;2;3)$, $B(3;0;1)$ et $C(2;1;2)$.
\begin{enumerate}[label=\alph*)]
\item Les points $A$, $B$ et $C$ sont-ils alignés ? Justifier.
\item Déterminer une équation paramétrique de la droite $(AB)$.
\item Le point $D(5;-2;-1)$ appartient-il à la droite $(AB)$ ?
\end{enumerate}

\rubrique{Vecteurs de l'espace et coordonnées}

\exo{}\diff{1}
Dans un repère orthonormé $(O;\vec{i},\vec{j},\vec{k})$, on donne les vecteurs $\vec{u}(2;-1;3)$ et $\vec{v}(-4;2;-6)$.
\begin{enumerate}[label=\alph*)]
\item Calculer les coordonnées de $\vec{u}+\vec{v}$ et de $2\vec{u}-3\vec{v}$.
\item Les vecteurs $\vec{u}$ et $\vec{v}$ sont-ils colinéaires ? Justifier.
\end{enumerate}

\exo{}\diff{2}
On considère les points $A(1;0;2)$, $B(3;2;1)$ et $C(0;1;4)$.
\begin{enumerate}[label=\alph*)]
\item Calculer les coordonnées des vecteurs $\overrightarrow{AB}$ et $\overrightarrow{AC}$.
\item Déterminer les coordonnées du point $D$ tel que $ABCD$ soit un parallélogramme.
\item Calculer les coordonnées du milieu $I$ de $[AB]$ et du milieu $J$ de $[CD]$. Que constate-t-on ?
\end{enumerate}

\exo{}\diff{2}
Dans un repère orthonormé, on donne $\vec{u}(1;2;3)$ et $\vec{v}(-2;1;0)$.
\begin{enumerate}[label=\alph*)]
\item Calculer le produit scalaire $\vec{u}\cdot\vec{v}$.
\item En déduire si $\vec{u}$ et $\vec{v}$ sont orthogonaux.
\item Déterminer un vecteur $\vec{w}$ non nul orthogonal à la fois à $\vec{u}$ et $\vec{v}$.
\end{enumerate}

\rubrique{Parallélisme et orthogonalité}

\exo{}\diff{2}
$ABCDEFGH$ est un cube. On note $I$ le centre de la face $ABCD$ et $J$ le centre de la face $EFGH$.
\begin{enumerate}[label=\alph*)]
\item Montrer que la droite $(IJ)$ est parallèle au plan $(ABF)$.
\item Montrer que la droite $(IJ)$ est orthogonale à la droite $(AC)$.
\end{enumerate}

\exo{}\diff{3}
Soit $ABCD$ un tétraèdre régulier (toutes les arêtes sont égales). On note $I$ le milieu de $[AB]$ et $J$ le milieu de $[CD]$.
\begin{enumerate}[label=\alph*)]
\item Montrer que les droites $(IJ)$ et $(AB)$ sont orthogonales.
\item Montrer que les droites $(IJ)$ et $(CD)$ sont orthogonales.
\item Que peut-on dire de la droite $(IJ)$ par rapport aux droites $(AB)$ et $(CD)$ ?
\end{enumerate}

\rubrique{Sections planes de solides}

\exo{}\diff{2}
$ABCDEFGH$ est un cube. On considère le plan $P$ passant par les points $I$, $J$ et $K$ où $I$ est le milieu de $[AB]$, $J$ le milieu de $[BC]$ et $K$ le milieu de $[BF]$.
\begin{enumerate}[label=\alph*)]
\item Tracer la section du cube par le plan $P$.
\item Quelle est la nature de cette section ?
\end{enumerate}

\exo{}\diff{3}
$SABCD$ est une pyramide à base carrée $ABCD$ de centre $O$, et de sommet $S$. On note $I$ le milieu de $[SA]$ et $J$ le milieu de $[SB]$.
\begin{enumerate}[label=\alph*)]
\item Montrer que la droite $(IJ)$ est parallèle au plan $(ABCD)$.
\item Déterminer l'intersection du plan $(IJK)$ où $K$ est le milieu de $[SC]$ avec la face $SCD$.
\item En déduire la section de la pyramide par le plan $(IJK)$.
\end{enumerate}

\section{Exercices d'entraînement}

\rubrique{Positions relatives}

\exo{}\diff{1}
Dans un cube $ABCDEFGH$, préciser les positions relatives :
\begin{enumerate}[label=\alph*)]
\item de la droite $(AB)$ et du plan $(CDH)$ ;
\item de la droite $(AG)$ et du plan $(BCF)$ ;
\item des droites $(AD)$ et $(FG)$ ;
\item des droites $(AH)$ et $(CF)$.
\end{enumerate}

\exo{}\diff{1}
Soit $ABCDEFGH$ un parallélépipède rectangle. Les droites $(AC)$ et $(EG)$ sont-elles coplanaires ? Justifier.

\exo{}\diff{2}
On considère un tétraèdre $ABCD$. $I$ est le milieu de $[AB]$, $J$ le milieu de $[AC]$ et $K$ le milieu de $[AD]$.
\begin{enumerate}[label=\alph*)]
\item Montrer que le plan $(IJK)$ est parallèle au plan $(BCD)$.
\item En déduire que $IJ = \frac{1}{2}BC$, $JK = \frac{1}{2}CD$ et $IK = \frac{1}{2}BD$.
\end{enumerate}

\exo{}\diff{2}
Dans un repère orthonormé, on donne $A(2;1;0)$, $B(1;3;2)$ et $C(0;5;4)$.
\begin{enumerate}[label=\alph*)]
\item Montrer que $A$, $B$ et $C$ sont alignés.
\item Déterminer le réel $k$ tel que $\overrightarrow{AB}=k\overrightarrow{AC}$.
\end{enumerate}

\exo{}\diff{2}
Soit $P$ le plan passant par $A(1;0;1)$ et de vecteurs directeurs $\vec{u}(1;2;1)$ et $\vec{v}(0;1;2)$. Le point $B(2;3;4)$ appartient-il à $P$ ?

\rubrique{Vecteurs et coordonnées}

\exo{}\diff{1}
Dans un repère orthonormé, on donne $\vec{u}(3;-2;1)$ et $\vec{v}(-1;4;2)$. Calculer :
\begin{enumerate}[label=\alph*)]
\item $\vec{u}+\vec{v}$ et $\vec{u}-\vec{v}$ ;
\item $2\vec{u}+3\vec{v}$ ;
\item $\vec{u}\cdot\vec{v}$.
\end{enumerate}

\exo{}\diff{2}
On considère les points $A(1;2;3)$, $B(4;5;6)$ et $C(7;8;9)$.
\begin{enumerate}[label=\alph*)]
\item Calculer les coordonnées du vecteur $\overrightarrow{AB}$.
\item Déterminer les coordonnées du point $D$ tel que $\overrightarrow{AD}=2\overrightarrow{AB}$.
\item Les points $A$, $B$ et $C$ sont-ils alignés ?
\end{enumerate}

\exo{}\diff{2}
Soit $\vec{u}(1;2;3)$ et $\vec{v}(-2;1;4)$. Déterminer un vecteur $\vec{w}$ tel que $\vec{w}$ soit orthogonal à $\vec{u}$ et à $\vec{v}$.

\exo{}\diff{3}
Dans un repère orthonormé, on donne $A(1;0;2)$, $B(3;2;1)$ et $C(0;1;4)$.
\begin{enumerate}[label=\alph*)]
\item Calculer les coordonnées du point $G$ tel que $\overrightarrow{GA}+\overrightarrow{GB}+\overrightarrow{GC}=\vec{0}$.
\item Vérifier que $G$ est le centre de gravité du triangle $ABC$.
\end{enumerate}

\rubrique{Parallélisme et orthogonalité}

\exo{}\diff{2}
$ABCDEFGH$ est un cube. Montrer que la droite $(EC)$ est orthogonale au plan $(ABD)$.

\exo{}\diff{2}
Soit $ABCD$ un tétraèdre. On suppose que $(AB)$ est orthogonale à $(CD)$ et que $(AC)$ est orthogonale à $(BD)$. Montrer que $(AD)$ est orthogonale à $(BC)$.

\exo{}\diff{3}
Dans un repère orthonormé, on considère les points $A(1;2;3)$, $B(3;4;5)$ et $C(2;1;4)$.
\begin{enumerate}[label=\alph*)]
\item Calculer les produits scalaires $\overrightarrow{AB}\cdot\overrightarrow{AC}$ et $\overrightarrow{BA}\cdot\overrightarrow{BC}$.
\item En déduire la nature du triangle $ABC$.
\end{enumerate}

\rubrique{Sections planes}

\exo{}\diff{2}
$ABCDEFGH$ est un cube. On considère le plan $P$ passant par les milieux des arêtes $[AB]$, $[AD]$ et $[AE]$. Tracer la section du cube par $P$ et déterminer sa nature.

\exo{}\diff{3}
$SABCD$ est une pyramide à base carrée $ABCD$. $I$ est un point de $[SA]$, $J$ un point de $[SB]$ et $K$ un point de $[SC]$. Construire la section de la pyramide par le plan $(IJK)$.

\exo{}\diff{3}
On considère un tétraèdre $ABCD$. $I$ est le milieu de $[AB]$, $J$ le milieu de $[AC]$ et $K$ le milieu de $[BD]$. Déterminer la section du tétraèdre par le plan $(IJK)$.

\section{Problèmes type Probatoire/Bac}

\sujetbac[Géométrie dans l'espace et sections planes]
$ABCDEFGH$ est un cube d'arête $4$ cm. On considère le repère orthonormé $(A;\vec{AB},\vec{AD},\vec{AE})$.
\begin{enumerate}[label=\arabic*)]
\item Donner les coordonnées de tous les sommets du cube.
\item Soit $I$ le point de coordonnées $(2;0;4)$ et $J$ le point de coordonnées $(4;2;0)$.
\begin{enumerate}[label=\alph*)]
\item Montrer que $I$ et $J$ appartiennent à des arêtes du cube. Préciser lesquelles.
\item Calculer les coordonnées du vecteur $\overrightarrow{IJ}$.
\end{enumerate}
\item Soit $K$ le milieu de $[CG]$.
\begin{enumerate}[label=\alph*)]
\item Donner les coordonnées de $K$.
\item Montrer que les points $I$, $J$ et $K$ ne sont pas alignés.
\end{enumerate}
\item On considère le plan $P$ passant par $I$, $J$ et $K$.
\begin{enumerate}[label=\alph*)]
\item Déterminer un vecteur normal au plan $P$.
\item En déduire une équation cartésienne du plan $P$.
\end{enumerate}
\item Déterminer la section du cube par le plan $P$. On précisera la nature de cette section.
\end{enumerate}

\sujetbac[Vecteurs et positions relatives dans un tétraèdre]
$ABCD$ est un tétraèdre. On note $I$ le milieu de $[AB]$, $J$ le milieu de $[CD]$, $K$ le milieu de $[BC]$ et $L$ le milieu de $[AD]$.
\begin{enumerate}[label=\arabic*)]
\item Faire une figure.
\item Montrer que $\overrightarrow{IJ} = \frac{1}{2}(\overrightarrow{AC} + \overrightarrow{BD})$.
\item Montrer que $\overrightarrow{KL} = \frac{1}{2}(\overrightarrow{AC} + \overrightarrow{BD})$.
\item Que peut-on dire du quadrilatère $IKJL$ ?
\item On suppose que $AB = CD$ et que $(AB) \perp (CD)$. Montrer que $IKJL$ est un losange.
\end{enumerate}

\section{Corrigés}

\corrige{de l'exercice 1}
\begin{enumerate}[label=\alph*)]
\item Dans le cube $ABCDEFGH$, le plan $(EFG)$ est le plan de la face supérieure. La droite $(AB)$ est une arête de la face inférieure. Comme $(AB)$ est parallèle à $(EF)$ et que $(EF)$ est contenue dans $(EFG)$, la droite $(AB)$ est parallèle au plan $(EFG)$.
\item Le plan $(EFH)$ est le plan de la face supérieure. La droite $(AC)$ est une diagonale de la face $ABCD$. Elle coupe le plan $(EFH)$ au point $G$ (car $(AC)$ et $(EG)$ sont sécantes en $G$). Donc $(AC)$ est sécante au plan $(EFH)$.
\item Les droites $(AB)$ et $(GH)$ sont parallèles car ce sont deux arêtes opposées du cube.
\item Les droites $(AC)$ et $(FH)$ sont parallèles car ce sont les diagonales de deux faces parallèles.
\end{enumerate}

\corrige{de l'exercice 2}
\begin{enumerate}[label=\alph*)]
\item Dans le triangle $ABC$, $I$ est sur $[AB]$ et $J$ sur $[AC]$. D'après le théorème de Thalès dans l'espace, les droites $(IJ)$ et $(BC)$ sont parallèles.
\item Les droites $(IJ)$ et $(BC)$ sont parallèles, donc le plan $(IJK)$ contient une droite parallèle à $(BC)$. De plus, $(IJ)$ et $(IK)$ sont sécantes en $I$. Par le théorème du parallélisme, le plan $(IJK)$ est parallèle au plan $(BCD)$.
\end{enumerate}

\corrige{de l'exercice 3}
\begin{enumerate}[label=\alph*)]
\item $\overrightarrow{AB}(2;-2;-2)$ et $\overrightarrow{AC}(1;-1;-1)$. On a $\overrightarrow{AB}=2\overrightarrow{AC}$, donc les vecteurs sont colinéaires. Les points $A$, $B$ et $C$ sont alignés.
\item Une équation paramétrique de $(AB)$ est : $\begin{cases} x=1+2t \\ y=2-2t \\ z=3-2t \end{cases}$, $t\in\R$.
\item Pour $D(5;-2;-1)$, on cherche $t$ tel que $\begin{cases} 5=1+2t \\ -2=2-2t \\ -1=3-2t \end{cases}$. La première donne $t=2$, la deuxième $t=2$, la troisième $t=2$. Donc $D$ appartient à $(AB)$.
\end{enumerate}

\corrige{de l'exercice 4}
\begin{enumerate}[label=\alph*)]
\item $\vec{u}+\vec{v}=(-2;1;-3)$ ; $2\vec{u}-3\vec{v}=(4+12;-2-6;6+18)=(16;-8;24)$.
\item $\vec{v}=-2\vec{u}$ car $(-4;2;-6)=-2(2;-1;3)$. Donc $\vec{u}$ et $\vec{v}$ sont colinéaires.
\end{enumerate}

\corrige{de l'exercice 5}
\begin{enumerate}[label=\alph*)]
\item $\overrightarrow{AB}(2;2;-1)$ ; $\overrightarrow{AC}(-1;1;2)$.
\item $D$ tel que $ABCD$ parallélogramme $\Leftrightarrow \overrightarrow{AB}=\overrightarrow{DC}$. Soit $D(x;y;z)$, alors $\overrightarrow{DC}(0-x;1-y;4-z)=(2;2;-1)$. D'où $x=-2$, $y=-1$, $z=5$. Donc $D(-2;-1;5)$.
\item $I$ milieu de $[AB]$ : $I(2;1;1,5)$. $J$ milieu de $[CD]$ : $J(-1;0;4,5)$. On constate que $I$ et $J$ sont différents, ce qui est normal car $ABCD$ n'est pas un parallélogramme centré.
\end{enumerate}

\corrige{de l'exercice 6}
\begin{enumerate}[label=\alph*)]
\item $\vec{u}\cdot\vec{v}=1\times(-2)+2\times1+3\times0=-2+2+0=0$.
\item $\vec{u}\cdot\vec{v}=0$, donc $\vec{u}$ et $\vec{v}$ sont orthogonaux.
\item Un vecteur $\vec{w}$ orthogonal à $\vec{u}$ et $\vec{v}$ peut être obtenu par le produit vectoriel (hors programme) ou par résolution. On cherche $\vec{w}(x;y;z)$ tel que $\vec{w}\cdot\vec{u}=0$ et $\vec{w}\cdot\vec{v}=0$. Soit $\begin{cases} x+2y+3z=0 \\ -2x+y=0 \end{cases}$. De la deuxième, $y=2x$. En substituant dans la première : $x+4x+3z=0 \Rightarrow 5x+3z=0 \Rightarrow z=-\frac{5}{3}x$. En prenant $x=3$, on a $y=6$, $z=-5$. Donc $\vec{w}(3;6;-5)$ convient.
\end{enumerate}

\corrige{de l'exercice 7}
\begin{enumerate}[label=\alph*)]
\item Dans le cube, $I$ est le centre de $ABCD$ et $J$ le centre de $EFGH$. La droite $(IJ)$ est parallèle à $(AE)$ (arête verticale). Or $(AE)$ est parallèle au plan $(ABF)$ car $(AE)$ est parallèle à $(BF)$ contenue dans $(ABF)$. Donc $(IJ)$ est parallèle à $(ABF)$.
\item $(AC)$ est une diagonale de la face $ABCD$. $(IJ)$ est verticale. Dans le plan horizontal $ABCD$, $(AC)$ est orthogonale à toute droite verticale, donc $(IJ)\perp(AC)$.
\end{enumerate}

\corrige{de l'exercice 8}
\begin{enumerate}[label=\alph*)]
\item Dans le tétraèdre régulier, toutes les arêtes sont égales. $I$ est le milieu de $[AB]$. Dans le triangle $AIB$, $IA=IB$ car $I$ est milieu. De plus, $AI=BI$ et $AB$ est une arête. Par symétrie, $(IJ)$ est la médiatrice de $[AB]$, donc $(IJ)\perp(AB)$.
\item De même, $J$ est le milieu de $[CD]$, et par symétrie, $(IJ)$ est la médiatrice de $[CD]$, donc $(IJ)\perp(CD)$.
\item La droite $(IJ)$ est orthogonale à la fois à $(AB)$ et $(CD)$. C'est la perpendiculaire commune aux deux droites $(AB)$ et $(CD)$.
\end{enumerate}

\corrige{de l'exercice 9}
\begin{enumerate}[label=\alph*)]
\item Dans le cube, $I$ milieu de $[AB]$, $J$ milieu de $[BC]$, $K$ milieu de $[BF]$. Le plan $(IJK)$ coupe la face $ABCD$ selon le segment $[IJ]$, la face $ABFE$ selon $[IK]$, la face $BCGF$ selon $[JK]$. La section est un triangle $IJK$.
\item $IJ$ est une diagonale d'un carré de côté moitié, $IK$ et $JK$ aussi. Donc $IJ=IK=JK$. La section est un triangle équilatéral.
\end{enumerate}

\corrige{de l'exercice 10}
\begin{enumerate}[label=\alph*)]
\item $I$ milieu de $[SA]$ et $J$ milieu de $[SB]$. Dans le triangle $SAB$, $(IJ)$ est une droite des milieux, donc $(IJ)\parallel(AB)$. Comme $(AB)$ est contenue dans le plan $(ABCD)$, $(IJ)$ est parallèle à ce plan.
\item $K$ milieu de $[SC]$. Dans le triangle $SBC$, $(JK)$ est une droite des milieux, donc $(JK)\parallel(BC)$. Dans le triangle $SCD$, la droite passant par $K$ parallèle à $(CD)$ coupe $[SD]$ en son milieu $L$. Donc le plan $(IJK)$ coupe la face $SCD$ selon $[KL]$.
\item La section de la pyramide par le plan $(IJK)$ est le quadrilatère $IJKL$. C'est un parallélogramme car $IJ\parallel AB\parallel CD\parallel KL$ et $JK\parallel BC\parallel AD\parallel IL$.
\end{enumerate}

\corrige{du problème 1}
\begin{enumerate}[label=\arabic*)]
\item Dans le repère $(A;\vec{AB},\vec{AD},\vec{AE})$, on a : $A(0;0;0)$, $B(4;0;0)$, $C(4;4;0)$, $D(0;4;0)$, $E(0;0;4)$, $F(4;0;4)$, $G(4;4;4)$, $H(0;4;4)$.
\item \begin{enumerate}[label=\alph*)]
\item $I(2;0;4)$ : $x=2$, $y=0$, $z=4$. $I$ est sur l'arête $[AE]$ car $A(0;0;4)$ et $E(0;0;4)$ ? Non, $I$ a $x=2$, donc il est sur $[EF]$ ? Vérifions : $E(0;0;4)$, $F(4;0;4)$, donc $I$ est sur $[EF]$. $J(4;2;0)$ : $J$ est sur $[BC]$ car $B(4;0;0)$ et $C(4;4;0)$.
\item $\overrightarrow{IJ}(4-2;2-0;0-4)=(2;2;-4)$.
\end{enumerate}
\item \begin{enumerate}[label=\alph*)]
\item $K$ milieu de $[CG]$ : $C(4;4;0)$, $G(4;4;4)$, donc $K(4;4;2)$.
\item $\overrightarrow{IK}(4-2;4-0;2-4)=(2;4;-2)$ et $\overrightarrow{IJ}(2;2;-4)$. Ces vecteurs ne sont pas colinéaires car $\frac{2}{2}\neq\frac{4}{2}$. Donc $I$, $J$, $K$ non alignés.
\end{enumerate}
\item \begin{enumerate}[label=\alph*)]
\item Un vecteur normal $\vec{n}$ au plan $P$ vérifie $\vec{n}\cdot\overrightarrow{IJ}=0$ et $\vec{n}\cdot\overrightarrow{IK}=0$. Soit $\vec{n}(a;b;c)$. On a $\begin{cases} 2a+2b-4c=0 \\ 2a+4b-2c=0 \end{cases}$. En soustrayant : $-2b-2c=0 \Rightarrow b=-c$. En substituant : $2a+2(-c)-4c=0 \Rightarrow 2a-6c=0 \Rightarrow a=3c$. En prenant $c=1$, on a $a=3$, $b=-1$. Donc $\vec{n}(3;-1;1)$.
\item Équation du plan $P$ : $3x-y+z+d=0$. Avec $I(2;0;4)$ : $6-0+4+d=0 \Rightarrow d=-10$. Donc $P:3x-y+z-10=0$.
\end{enumerate}
\item La section du cube par $P$ est un hexagone. Les points d'intersection de $P$ avec les arêtes sont : sur $[EF]$ : $I$ ; sur $[BC]$ : $J$ ; sur $[CG]$ : $K$ ; sur $[DH]$ : $L$ ; sur $[AE]$ : $M$ ; sur $[AB]$ : $N$. La section est un hexagone régulier.
\end{enumerate}

\corrige{du problème 2}
\begin{enumerate}[label=\arabic*)]
\item Figure à faire.
\item $\overrightarrow{IJ} = \overrightarrow{IA} + \overrightarrow{AJ} = \frac{1}{2}\overrightarrow{BA} + \frac{1}{2}\overrightarrow{AC} + \frac{1}{2}\overrightarrow{AD} = \frac{1}{2}(\overrightarrow{BA}+\overrightarrow{AC}+\overrightarrow{AD}) = \frac{1}{2}(\overrightarrow{BC}+\overrightarrow{AD})$. Mais $\overrightarrow{BC} = \overrightarrow{BA}+\overrightarrow{AC}$ et $\overrightarrow{AD} = \overrightarrow{AB}+\overrightarrow{BD}$. En sommant : $\overrightarrow{BC}+\overrightarrow{AD} = \overrightarrow{BA}+\overrightarrow{AC}+\overrightarrow{AB}+\overrightarrow{BD} = \overrightarrow{AC}+\overrightarrow{BD}$. Donc $\overrightarrow{IJ} = \frac{1}{2}(\overrightarrow{AC}+\overrightarrow{BD})$.
\item $\overrightarrow{KL} = \overrightarrow{KB}+\overrightarrow{BL} = \frac{1}{2}\overrightarrow{CB} + \frac{1}{2}\overrightarrow{BA} = \frac{1}{2}(\overrightarrow{CB}+\overrightarrow{BA}) = \frac{1}{2}\overrightarrow{CA} = -\frac{1}{2}\overrightarrow{AC}$. Ce n'est pas égal. Reprenons : $K$ milieu de $[BC]$, $L$ milieu de $[AD]$. $\overrightarrow{KL} = \overrightarrow{KA}+\overrightarrow{AL} = \frac{1}{2}(\overrightarrow{BA}+\overrightarrow{CA}) + \frac{1}{2}\overrightarrow{AD} = \frac{1}{2}(\overrightarrow{BA}+\overrightarrow{CA}+\overrightarrow{AD}) = \frac{1}{2}(\overrightarrow{BD}+\overrightarrow{CD})$. En utilisant $\overrightarrow{BD} = \overrightarrow{BA}+\overrightarrow{AD}$ et $\overrightarrow{CD} = \overrightarrow{CA}+\overrightarrow{AD}$, on obtient $\overrightarrow{KL} = \frac{1}{2}(\overrightarrow{BA}+\overrightarrow{AD}+\overrightarrow{CA}+\overrightarrow{AD}) = \frac{1}{2}(\overrightarrow{BA}+\overrightarrow{CA}+2\overrightarrow{AD})$. Ce n'est pas égal à $\frac{1}{2}(\overrightarrow{AC}+\overrightarrow{BD})$. Il y a une erreur. En fait, $\overrightarrow{KL} = \overrightarrow{KA}+\overrightarrow{AL} = \frac{1}{2}\overrightarrow{BA} + \frac{1}{2}\overrightarrow{AD} = \frac{1}{2}(\overrightarrow{BA}+\overrightarrow{AD}) = \frac{1}{2}\overrightarrow{BD}$. Et $\overrightarrow{IJ} = \frac{1}{2}(\overrightarrow{AC}+\overrightarrow{BD})$. Donc $\overrightarrow{KL} \neq \overrightarrow{IJ}$ en général. Le quadrilatère $IKJL$ n'est pas un parallélogramme.
\item $I$, $K$, $J$, $L$ sont les milieux des arêtes. $IK$ est une droite des milieux dans $ABC$, donc $IK\parallel AC$ et $IK=\frac{1}{2}AC$. $JL$ est une droite des milieux dans $ACD$, donc $JL\parallel AC$ et $JL=\frac{1}{2}AC$. Donc $IK\parallel JL$ et $IK=JL$. De même, $IL\parallel BD$ et $IL=\frac{1}{2}BD$, $KJ\parallel BD$ et $KJ=\frac{1}{2}BD$. Donc $IL\parallel KJ$ et $IL=KJ$. Ainsi $IKJL$ est un parallélogramme.
\item Si $AB=CD$ et $(AB)\perp(CD)$, alors $AC=BD$ ? Non. Dans un tétraèdre, $AB=CD$ n'implique pas $AC=BD$. Mais $IK=\frac{1}{2}AC$ et $IL=\frac{1}{2}BD$. Si $AC=BD$, alors $IK=IL$ et le parallélogramme $IKJL$ est un losange. L'énoncé dit $AB=CD$, ce qui est différent. Il faut vérifier : $IKJL$ est un losange si $IK=IL$, soit $AC=BD$. Or $AB=CD$ et $(AB)\perp(CD)$ ne donne pas $AC=BD$. Il y a une erreur dans l'énoncé. On peut montrer que si $AB=CD$ et $(AB)\perp(CD)$, alors $AC=BD$ ? Non. Donc le problème est mal posé. On corrige : On suppose $AC=BD$ et $(AC)\perp(BD)$. Alors $IK=IL$ et $IKJL$ est un losange.
\end{enumerate}